Section~\ref{sec:ablation} ablates components \emph{within} the prompt while always supplying profiling
data. The prior question---whether the profiling data helps at all---requires an arm with no profile
data, which we add here.

\paragraph{Arms.}
The \texttt{pcrl} arm is the standard prompt carrying PC-sampling and roofline summaries only (no
derived kernel-budget, occupancy or diagnostic summaries). The \texttt{code\_only} arm removes the
entire performance-analysis block. Removing it alone would bias the comparison, because four
instruction sentences refer to the data---most importantly the rule \textit{``if no safe substantive
optimization can be justified from the provided performance data, keep the code semantically
unchanged''}, which would invite a no-op when no data is present. Those four sentences are therefore
reworded to refer to the model's analysis of the source; the two prompts are otherwise byte-identical
(verified: the diff is the removed block plus those four sentences). Both arms were regenerated with
the same code base, and all pairs use the same correctness gates and interleaved measurement as
Section~\ref{sec:realapps}.

\paragraph{Scoring.}
HeCBench pairs are measured on a shared H100 node where per-repetition spread is sometimes larger than
any plausible effect, so we score a pair as a \emph{tie} when $|\mathrm{pcrl}/\mathrm{co}-1|$ is smaller
than the largest repetition spread measured for that pair---that is, when the difference between the
arms does not exceed its own measurement noise. Two applications (\texttt{boxfilter},
\texttt{ccsd-trpdrv}) are excluded because the \emph{unmodified} program exceeds our $600$\,s per-run
cap on that node, and one (\texttt{background-subtract}, GPT-5) because the \texttt{code\_only} arm
produced no correct candidate at all, which we report under correctness rather than speed. One
caveat carries through the whole benchmark half of this comparison: on this node the benchmarks' own
verification gives no usable verdict for the \emph{unmodified} program---five of nine report FAIL and
the rest report nothing---so a candidate's verdict is informative only where it differs from the
pristine one (see Limitations).

\paragraph{Result.}
Profiling data helps on the large real application and makes no measurable difference on the small
benchmarks (Table~\ref{tab:profabl}). On Cholla both models are clearly faster with the profile
($+11.5\%$ for GPT-5, $+7.7\%$ for Gemini), and a second, independently generated \texttt{pcrl} sample
keeps both on the same side of the comparison (GPT-5 $1.187$, Gemini $1.023$ relative to
\texttt{code\_only}), although the Gemini replicate is much weaker than its first sample. On
AutoDock-GPU both arms reach the same ${\sim}1.03\times$. All seventeen scored HeCBench pairs are ties
by the rule above: the largest arm-to-arm difference is $7.8\%$ for GPT-5 (\texttt{adv}, inside a
$31.8\%$ spread) and $33.2\%$ for Gemini (\texttt{background-subtract}, inside a $147.6\%$ spread), and
the eight pairs whose spread is under $5\%$ all land between $0.998$ and $1.003$. Not one benchmark
shows a difference that its own measurement noise does not already cover.


\begin{table}[t]
\centering
\footnotesize
\caption{Profile-data ablation: speedup against the unmodified application for the prompt without
  profile data (\texttt{co}) and with PC sampling + roofline (\texttt{pcrl}), and their ratio.
  A HeCBench pair is a tie when the arm-to-arm difference is smaller than that pair's repetition
  spread; the entry gives the largest difference observed and the spread it sits inside.}
\label{tab:profabl}
\setlength{\tabcolsep}{2pt}
\begin{tabular}{@{}llrrr@{}}
\toprule
\textbf{Application} & \textbf{Model} & \textbf{co} & \textbf{pcrl} & \textbf{pcrl/co} \\
\midrule
\multirow{2}{*}{Cholla}
  & GPT-5  & 1.324$\times$ & \textbf{1.476}$\times$ & \textbf{1.115} \\
  & Gemini & 1.471$\times$ & \textbf{1.584}$\times$ & \textbf{1.077} \\
\midrule
\multirow{2}{*}{AutoDock-GPU}
  & GPT-5  & 1.030$\times$ & 1.031$\times$ & 1.001 \\
  & Gemini & 1.035$\times$ & 1.035$\times$ & 1.000 \\
\midrule
HeCBench (8 pairs) & GPT-5  & \multicolumn{3}{c}{all ties} \\
HeCBench (9 pairs) & Gemini & \multicolumn{3}{c}{all ties} \\
\bottomrule
\end{tabular}
\end{table}

\paragraph{Correctness, not only speed.}
Two HeCBench pairs favour the profile arm on correctness rather than time. On
\texttt{background-subtract}, GPT-5 without profile data failed the application's own output check in
all three repair iterations and produced no usable candidate, while the \texttt{pcrl} arm produced a
correct one; on \texttt{atomicSystemWide}, the \texttt{code\_only} candidate fails verification and the
\texttt{pcrl} candidate does not.

\paragraph{Interpretation.}
The pattern is consistent with the workloads rather than with the models. A HeCBench kernel is a few
hundred lines with one obvious hot loop, and the profile tells a model little it cannot infer from the
source. Cholla's HLLC solver is 307 lines inside a 21k-line CUDA application whose cost is spread over
several kernels; there the profile identifies both the kernel worth editing and the reason it is slow
(latency-bound at $19.7\%$ DRAM and $40.4\%$ SM throughput, occupancy $23\%$ limited by 109 registers
per thread). We therefore expect profile data to pay off in proportion to how hard the bottleneck is to
locate from source alone---which is precisely the regime real applications occupy.
A separate question is what happens when that data is not merely supplied but \emph{interpreted} into
named remedies before the model sees it; Section~\ref{sec:guidvar} measures that arm.
